\FloatBarrier
\newpage
\section{数据的获取}
\label{sec:prob-sampling}

给定模型可以研究可能出现的数据。实际统计往往从另一端开始：面对尚不完全了解的总体，怎样获得有用的数据？抽样方式决定哪些对象可能入选，也决定随后概率计算与统计判断的依据。

\subsection{总体与样本}

\begin{definition}
研究所涉及的全部对象或其测量值组成{\heiti 总体}；实际选取并观察的一部分对象或其测量值组成{\heiti 样本}，观测个数称为{\heiti 样本量}。
\end{definition}

研究一批产品的不合格比例，总体是这批产品，记录变量为是否不合格。研究某种装置的测量误差，总体也可以由在固定条件下反复测量的概率分布描述。前者通常有一个明确的有限对象集合，后者主要由试验机制与分布模型确定。

总体比例、总体均值等描述总体的量称为{\heiti 参数}。样本中的比例、均值是由实际数据计算的量。取样以前，它们随样本变化；取样以后，它们成为具体的观测值。后面的抽样分布将进一步研究这种区别。

明确总体还包括明确时间、范围和测量方法。例如调查某校学生的每周锻炼时间，需要确定研究哪一学期、怎样计算锻炼时间，以及缺少记录时怎样处理。不同的记录规则可能形成不同的数据。

\subsection{简单随机抽样}

\begin{definition}
从\(N\)个不同对象中不放回抽取\(n\)个，若每个\(n\)元子集都具有概率\(1/\binom Nn\)，称为{\heiti 简单随机抽样}。
\end{definition}

可以将对象编号，用随机顺序取前\(n\)个，或均匀抽取不同编号。每个对象的入选概率为
\[
 \frac{\binom{N-1}{n-1}}{\binom Nn}=\frac nN.
\]
但仅有单个对象入选概率相同，还不足以保证所有子集等可能。例如四人中只在\(\{1,2\}\)、\(\{3,4\}\)两组之间等可能选择，每人的入选概率都是\(1/2\)，却没有给其他二人组任何机会。

有限总体中不放回抽样的各次结果相互关联。若采用有放回的独立抽取，或直接假定观测来自同一模型的独立重复试验，就得到另一种抽样模型；使用均值方差等公式时应按实际机制选取。

\subsection{分层随机抽样}

若总体由不同组别组成，可以先按有意义的特征分层，再在各层内随机抽取。设第\(h\)层有\(N_h\)个对象，样本量为\(n_h\)，总体量为\(N=\sum_hN_h\)。按比例分配时，令\(n_h/n\)尽量接近\(N_h/N\)。

若各层样本均值为\(\bar x_h\)，估计全体平均时应使用总体权重：
\begin{equation}\label{eq:prob-stratified-estimate}
 \widehat\mu=\sum_h\frac{N_h}{N}\bar x_h.
\end{equation}
比例分配时，它与全部样本直接计算的均值一致；非比例分配时，样本中的各层比例已经改变，需用上式恢复总体的构成。

例如两个年级分别有900人、100人，各抽20人，样本平均得分分别为60、90。按总体构成估计的平均为\(0.9\cdot60+0.1\cdot90=63\)，直接合并40人的平均则为75。后一个数描述当前样本的构成，前一个数用于恢复原总体的权重。

分层使各类对象有计划地进入样本，是否能降低估计的波动，还与层内差异和分配方式有关。本节先建立抽样机制与加权依据。

\subsection{研究设计}

随机抽样用于从目标总体取得样本；随机分组用于把参加者分配给不同处理。两者承担不同任务。随机抽样支持从样本向总体的推断，随机分组则帮助比较处理效果，减少原有差异对比较的系统影响。

观察性研究记录原有状态，例如比较自愿参加训练与未参加训练者的成绩。两组还可能在基础、时间投入等方面不同。对照试验则由研究者安排处理，在条件允许时随机分组，并尽量保持其他条件可比。即使观察到明显相关，因果解释仍需要研究设计提供依据。

自愿回应、只调查容易接触的人、部分对象未回应，以及测量口径不一致，都可能带来系统偏差。扩大同一种有偏样本，主要增加同一获取方式下的记录，不能自行恢复原来缺失的总体信息。

\begin{example}
某校准备调查全校学生的睡眠时间。比较以下方案：在图书馆晚间离馆者中调查；在全校学生名单中简单随机抽取；按年级分层随机抽取。说明各方案的目标总体、可能遗漏的对象，以及分层比例不同时怎样汇总平均值。
\end{example}

\begin{exercise}
\mbox{}\par\nobreak

（1）某总体分三层，人数为500、300、200，计划抽取100人。给出比例分配方案。若改为各层抽30人，三个样本均值为12、15、18，求总体均值的分层估计。

（2）为了研究某训练方法，从报名者中随机分为训练组与对照组。说明随机分组解决了什么问题；若希望结论适用于所有高中生，还需要考虑哪些取样信息？

（3）从编号1至20的对象中随机选一个起点\(1,\ldots,5\)，再每隔5个取一个。求每个对象的入选概率，并判断这是否是样本量4的简单随机抽样。

（4）（据2022年新高考II卷第19题改编）随机调查某地区100名某疾病患者，将年龄分入\([0,10),[10,20),\ldots,[80,90]\)，得到下表：
\[
\begin{array}{c|rrrrrrrrr}
 \text{组中值（岁）}&5&15&25&35&45&55&65&75&85\\\hline
 \text{频数}&1&2&12&17&23&20&17&6&2
\end{array}
\]
各组以组中值代表，估计患者的平均年龄及年龄在\([20,70)\)内的概率。若全地区患病率为\(0.1\%\)，年龄在\([40,50)\)内的人口占\(16\%\)，估计从这一年龄段随机选一人时其患病的概率，精确到\(0.0001\)。说明最后一问所用的总体与患者样本之间的关系。
\end{exercise}
